\documentclass[11pt,a4paper]{article}
\usepackage[margin=25mm, headheight=14pt]{geometry}
\input{preamble}
\usepackage[numbers,sort&compress]{natbib}
\bibliographystyle{plainnat}

\title{\vspace{-1.5cm}\bfseries The TYTFS 2024 study files as a PyPSA network\\[2mm]
\large\mdseries\color{inksoft} What was built, what it was checked against,
and what it is not good for}
\author{\normalsize\color{inksoft} Built with Claude Code}
\date{\normalsize\color{inksoft}\today}

\begin{document}
\maketitle
\thispagestyle{empty}

\begin{abstract}
\noindent
EirGrid publishes the Ten Year Transmission Forecast Statement (TYTFS) as four
PSS/E load-flow cases. This work turns them into PyPSA networks, puts
coordinates on the buses from OpenStreetMap, extracts the North-West region and
reconciles it against a hand-built dataset, attaches a week of hourly profiles,
and packages the result as a participant kit with six worked examples.

Everything below is reproducible from the repository. Every number quoted is
produced by code in it, and the sections say which function. Where something
could not be done -- the real weather data is behind a blocked host -- that is
stated rather than papered over.

This revision is the result of a deliberate red-teaming pass over the first
draft. Every external claim now carries a citation, every measured number was
re-derived from the data, and Section~\ref{sec:corrections} lists the four
things the first draft got wrong. One of them was a limitation that turned out
to be false in the helpful direction.
\end{abstract}

\vspace{4mm}
\noindent\colorbox{surface}{%
\begin{minipage}{\dimexpr\textwidth-2\fboxsep}
\small
\textbf{The short version.}
\begin{itemize}[leftmargin=1.2em]
\item Four TYTFS cases $\rightarrow$ eight PyPSA networks (two voltage scopes
      each). All connected, all solve a DC power flow, all solve a linear
      optimal power flow with HiGHS.
\item The conversion is checked against the cases' \emph{own} solved voltage
      angles: correlation $\geq 0.99$ on seven of eight networks.
\item 87\% of transmission buses geocoded; cross-checked against EirGrid's own
      station register at a median separation of \textbf{14 metres}.
\item Two independent methods -- a DC flow on the 2024 network and a full LOPF
      on the 2033 network with synthetic weather -- pick out the \emph{same}
      binding circuit: Letterkenny--Strabane, 93~MVA, into Northern Ireland.
\item The stated line lengths are \emph{real}: they correlate at $r = 0.983$
      with the great-circle distance between independently geocoded endpoints.
      The first draft of this report called them placeholders. It was wrong.
\item 215 automated tests pass. Two of them check the linear algebra against
      things that do not share code with it.
\end{itemize}
\end{minipage}}

\clearpage
\tableofcontents
\clearpage

% ===================================================================
\section{What this is, and how to run it}

The repository has two halves.

\textbf{The build half} reads the PSS/E files and produces everything else. It
is a handful of single-purpose modules with a command-line interface each:

\begin{lstlisting}[language=shell]
python geocode.py fetch          # OpenStreetMap -> data/raw/
python geocode.py match data/TYTFS2024_studyfiles/*_V35.raw
python pypsa_net.py build        # -> data/pypsa/<case>_<scope>/
python pypsa_net.py verify       # connectivity, DC power flow, LOPF
python northwest.py extract      # the North-West region, both views
python synthetic.py build        # a year of profiles, seeded
python build_kit.py              # -> participant-kit/
python -m pytest                 # 215 tests
\end{lstlisting}

\textbf{The participant half} is \file{participant-kit/}, a self-contained
directory that can be lifted out whole. Nothing in it imports anything from the
build half. It ships the networks, a loader, the flow linear algebra, and six
examples:

\begin{lstlisting}[language=shell]
python -m venv .venv && source .venv/bin/activate
pip install -r requirements.txt
python examples/a_dc_power_flow.py
\end{lstlisting}

That is the whole quickstart. Measured from a clean virtual environment it
takes 42 seconds and ends in a plotted map.

\subsection{The source data}

Four cases, each a PSS/E v35 raw file:

\begin{table}[H]\centering\small
\begin{tabular}{llrrr}
\toprule
\textbf{case} & \textbf{condition} & \textbf{buses} & \textbf{peak demand} &
\textbf{capacity} \\
\midrule
WP2024 & winter peak, 2024 network   & 646 & 7{,}325 MW & 22{,}650 MW \\
SV2024 & summer valley, 2024 network & 639 & 4{,}948 MW & 14{,}591 MW \\
WP2033 & winter peak, 2033 network   & 754 & 8{,}792 MW & 42{,}574 MW \\
SV2033 & summer valley, 2033 network & 754 & 6{,}058 MW & 18{,}283 MW \\
\bottomrule
\end{tabular}
\caption{The four TYTFS cases as they appear in the transmission-scope
networks. WP2033 is the interesting one: 42.6 GW of connected plant against an
8.8 GW peak. That headroom is where constraint and curtailment live.}
\end{table}

A snapshot is not a time series. Each case is one hour, with a solved dispatch
and solved voltages. Section~\ref{sec:profiles} is about what to do with that.

\where{psse.py}{read\_raw(), read\_all(), generators(), loads(), ac\_branches()}

% ===================================================================
\section{Turning a PSS/E case into a PyPSA network}
\label{sec:conversion}

Six decisions matter, in the sense that getting one wrong produces a network
that still solves and whose every downstream number is wrong. Each is below
with the measurement it was made on.

\subsection{The per-unit base: 100 MVA, read rather than assumed}

Every case opens with \code{SBASE = 100.00} and \code{BASFRQ = 50}
\citep{eirgrid_tytfs2024}. The build asserts it and refuses a case on any other
base -- a case on 200 MVA would
produce a network that solves and whose every impedance is out by exactly two.

PSS/E gives impedances per-unit on that base and on the bus's kV. PyPSA wants
ohms for a line and per-unit on the transformer's own rating for a transformer,
so there are two conversions:

\begin{table}[H]\centering\small
\begin{tabular}{lll}
\toprule
& \textbf{PyPSA units} & \textbf{conversion} \\
\midrule
\code{Line.r}, \code{Line.x}        & ohms & $R_{pu}\times \mathrm{kV}^2/100$ \\
\code{Line.b}                        & siemens & $B_{pu}\times 100/\mathrm{kV}^2$ \\
\code{Transformer.r}, \code{.x}      & pu on its own \code{s\_nom} & $X_{pu}\times \mathrm{s\_nom}/100$ \\
\bottomrule
\end{tabular}
\end{table}

\where{pypsa\_net.py}{\_ohms(), \_siemens(), \_z\_on\_system\_base()}

\subsection{\texorpdfstring{\code{RATE1} is taken as the rating -- by inference}{RATE1 is taken as the rating - by inference}}
\label{sec:rate}

RATE1, RATE2 and RATE3 are populated; RATE4--RATE12 are zero everywhere.

\textbf{A caveat this report owes the reader.} PSS/E v35 carries up to twelve
rating sets per branch, and Siemens' own manual -- which would say what, if
anything, the format fixes about their meaning -- is not publicly available.
Open-source readers of the format document the older three-rating
\code{RATEA}/\code{RATEB}/\code{RATEC} layout without settling the semantics of
the twelve-rating one \citep{powsybl_psse}. So what follows is an
\emph{inference from the data}, not a documented property of the format, and
the first draft of this report stated it as though it were the latter.

Two facts about the data, measured across all four cases:

\begin{itemize}
\item \code{RATE1 == RATE2} in \textbf{every one of the 4{,}841 branch records
      across the four cases} (1{,}129 in WP2024, 1{,}120 in SV2024, 1{,}296 in
      each 2033 case). Not approximately -- exactly.
\item \code{RATE3/RATE1} is 1.0 or 1.1 in the two winter cases. In the two
      summer cases 97 records come out at 1.10021 instead, which is a rating
      that was rounded before it was written rather than a third convention.
\end{itemize}

The natural reading is the normal / long-term-emergency / short-term-emergency
convention, with the LTE set equal to the normal rating and the STE at 110\%.
\code{s\_nom} is \code{RATE1}, the most conservative of the three. If that
reading is wrong, the direction of the error is known: \code{s\_nom} would be
too small, and every loading figure here too high.

\where{pypsa\_net.py}{CONTINUOUS\_RATE; test\_pypsa\_net.py $\rightarrow$ test\_rate2\_equals\_rate1}

\subsection{The 18 unrated branches: what they actually are}

Eighteen AC branches at 110~kV and above carry no rating in WP2024 -- nine at
\code{RATE1 = 9999} and nine at \code{RATE1 = 0}. Zero is the more dangerous of
the two, because it looks like a number: put into \code{s\_nom} it makes the
element unusable and the optimisation infeasible.

Looking at what they are settles the treatment. All eighteen have $R = 0$,
$|X| \leq 0.01$, $B = 0$ and $\mathrm{LEN} = 0$ -- as do all twenty in each
2033 case, where the count is higher. Every one is a station coupler: a busbar section, a capacitor or reactor stub, an SVC, a generator
terminal tie. Not one is a circuit whose rating was forgotten.

\textbf{The rule, in two parts.} A coupler is bounded by what is attached to
it, so \code{s\_nom} is the smaller of the two ends' summed rated incident
capacity -- a real limit, not a large number. Everything else keeps the file's
9999.

That second half was arrived at by trying the other rule. Bounding the
sub-110~kV distribution transformers the same way produced limits that bind on
the case's \emph{own solved flows} on 13 transformers, the worst at 156\% of
the bound it had been given, and made the full network's LOPF infeasible. There
is no honest way to infer a rating for a 110/38~kV transformer from the feeders
leaving the 38~kV bus, because load taps off that bus directly and the feeders
bound nothing.

\where{pypsa\_net.py}{\_is\_coupler(), \_fill\_unrated(), rating\_exceptions()}

\subsection{The sub-110~kV tail: a reduction, not a filter}

The networks come in two scopes, chosen by a \code{min\_kv} flag. Deleting the
sub-110~kV buses would disconnect the generation hanging off them, so instead
each removed bus is \emph{aggregated} into the retained bus it reaches by the
least-reactance path, and its generators and loads move with it. A wind farm at
690~V lands at the station it is actually connected to.

\begin{lstlisting}
model = pypsa_net.build(case, min_kv=pypsa_net.TRANSMISSION_KV)  # 110 kV+
model = pypsa_net.build(case, min_kv=0.0)                        # everything
\end{lstlisting}

\where{pypsa\_net.py}{\_retained(), \_aggregate\_to(), build()}

\subsection{The DC links}

Moyle, EWIC and Greenlink become PyPSA \code{Link} components -- a controllable
injection, not part of the linear AC network. Two of them have a far terminal
created for them because the file does not model the GB side. All three are
real HVDC links: Moyle (Scotland to Co.\ Antrim, in service 2001), the
East-West Interconnector (Portan in Co.\ Meath to Shotton in Wales, 500 MW, in
service 2012) and Greenlink (Great Island in Co.\ Wexford to Pembroke, in
commercial operation from 2025) \citep{eirgrid_interconnection}. The file names
the Greenlink terminal \code{GREENLINK}; OpenStreetMap names the converter
station after the townland, Campile, and it sits 400~m from the Great Island
220~kV substation -- which is why \file{geocode.py} carries that as a
hand-written alias with its reasoning attached.

\where{pypsa\_net.py}{INTERCONNECTORS, \_add\_links(), \_moyle\_efficiency()}

\subsection{The reference bus, and the bug that found it}
\label{sec:slack}

The case names its own: five buses carry \code{IDE = 3}, the PSS/E bus type
code for a swing bus \citep{powsybl_psse}. The rule is to honour
them where they survive the voltage floor, hand the role to whatever absorbed a
swing bus that was aggregated away, and otherwise give it to the bus carrying
the most in-service generation.

\textbf{This one is not cosmetic.} Determining the network topology makes PyPSA
choose a slack of its own -- the first generator in each sub-network -- and in
WP2033 that is a 21~MW hydro unit at Ardnacrusha behind its own step-up
transformer. A lossless DC flow has nowhere to put the case's 198~MW
generation-minus-demand gap except the reference, so that one machine absorbed
the lot: 856~MW through a 21~MW step-up, 172$^\circ$ across it, and the whole
network's angles thrown out by up to 100$^\circ$.

The fix is to assign the reference as a whole column \emph{after} the topology
has been determined, rather than label by label:

\begin{lstlisting}
# Assigning per-label leaves PyPSA's own choice in place.  Build the whole
# Series and set it as a column, after determine_network_topology().
control = pd.Series("PQ", index=n.generators.index)
control[chosen] = "Slack"
n.generators["control"] = control.values
\end{lstlisting}

Angle correlation against the case's own solved angles: \textbf{0.90 before,
0.99 after}.

\where{pypsa\_net.py}{\_slack\_buses(), build()}

\subsection{Verification}
\label{sec:verify}

\code{python pypsa\_net.py verify} runs all four cases at both scopes. The
network object and its linear power flow are PyPSA's \citep{pypsa}; the
optimisation is solved with HiGHS \citep{highs}.

\begin{table}[H]\centering\small
\begin{tabular}{llccrrl}
\toprule
\textbf{case} & \textbf{scope} & \textbf{connected} & \textbf{DC PF} &
\textbf{angle corr} & \textbf{angle sd} & \textbf{LOPF} \\
\midrule
WP2024 & transmission & yes & solved & 0.9961 & 0.93$^\circ$ & optimal \\
WP2024 & full         & yes & solved & 0.9915 & 1.19$^\circ$ & optimal \\
SV2024 & transmission & yes & solved & 0.9987 & 0.36$^\circ$ & optimal \\
SV2024 & full         & yes & solved & 0.9963 & 0.61$^\circ$ & optimal \\
WP2033 & transmission & yes & solved & 0.9907 & 1.94$^\circ$ & optimal \\
WP2033 & full         & yes & solved & 0.9048 & 4.35$^\circ$ & optimal \\
SV2033 & transmission & yes & solved & 0.9977 & 0.25$^\circ$ & optimal \\
SV2033 & full         & yes & solved & 0.9939 & 0.40$^\circ$ & optimal \\
\bottomrule
\end{tabular}
\caption{All eight networks. The angle check is the real verification: the raw
file records the solved voltage angle at every bus, so the conversion can be
checked against the case rather than against itself. If the reactances, the
three-winding star equivalents or the phase-shift signs were wrong, a linear
flow's angles would not track them.}
\end{table}

\where{pypsa\_net.py}{verify(), angle\_check(), \_connected\_with\_links()}

% ===================================================================
\section{Putting the buses on a map}
\label{sec:geocode}

Bus records have no coordinates and the only thing to join on is a
twelve-character name. \file{geocode.py} downloads every named
\code{power=substation} and \code{power=plant} on the island from Overpass
\citep{osm} -- 1{,}705 named objects, of which 1{,}270 carry a voltage tag and
367 are tagged at 110~kV or above -- and matches them against station names.

\textbf{Nothing is interpolated, averaged or guessed.} A station that cannot be
matched gets no coordinate and is reported by name with the candidate it
rejected and why.

\begin{figure}[H]\centering
\includegraphics[width=\textwidth]{g_geocoding.png}
\caption{Left: where the WP2024 transmission buses landed, coloured by how each
was matched. Right: the cross-check against the station layer of EirGrid's own
public Transmission Development Plan feature service
\citep{eirgrid_tdp2024_service}; 134 of the matched stations appear in both.
The median disagreement is \textbf{14~metres} -- the difference between one
corner of a substation and another -- and the 95th percentile is 252~metres.}
\label{fig:geocode}
\end{figure}

\begin{table}[H]\centering\small
\begin{tabular}{lrl}
\toprule
\textbf{method} & \textbf{buses} & \textbf{what it means} \\
\midrule
\code{exact}      & 309 & normalised names equal, voltage tag agrees \\
\code{coupled}    &  53 & joined to a placed bus by a coupler or its own transformer \\
\code{ni-code}    &  60 & Northern Ireland contraction resolved by subsequence \\
\code{truncated}  &  19 & the twelve-character name is a prefix of exactly one candidate \\
\code{fuzzy}      &  13 & token-sort score $\geq 92$ with $\geq 6$ points over the runner-up \\
\code{prefix}     &  10 & $\geq 6$-character prefix of exactly one candidate \\
\code{alias}      &   6 & resolved by hand, with the reasoning recorded \\
\code{ni-site}    &   6 & the same NIE site at its other voltage class \\
\code{exact-name} &   3 & names equal, candidate carries no voltage tag \\
\midrule
\textbf{failed}   & \textbf{70} & every one named, with its rejected candidate \\
\bottomrule
\end{tabular}
\caption{479 of 549 buses placed (87\%). Voltage is corroboration rather than a
key: OSM's voltage tag is present on 367 of the island's named substations and
absent on the rest, so agreement raises confidence, absence does not lower it,
and a contradiction rejects the match.}
\end{table}

\subsection{The three hard cases}

\textbf{Cathaleen's Fall} is spelled \emph{Cathleen's Fall} by OSM and appears
in TYTFS as \code{CATH FALL}. Resolved by an alias, recorded in the code with
its reasoning rather than as a magic string.

\textbf{Clogher} has four separate 110~kV buses -- split busbars. They are one
place, so \code{station\_of()} groups them before matching and they share one
coordinate.

\textbf{Northern Ireland} uses contractions rather than names: \code{DRUM1},
\code{GLEN1}. The trailing digit is a voltage class (1~=~110~kV, 2~=~275~kV) and
the letters are a subsequence of the real name. Two of these are genuinely
ambiguous -- \code{DRUM1} sits equally tightly in \emph{Drumnakelly} and
\emph{Drumlins Park Wind Farm} -- and are reported as failures rather than
guessed.

\where{geocode.py}{normalise(), \_subsequence\_span(), ALIASES, match(), crosscheck()}

\subsection{The one disagreement worth reporting}

\emph{A caveat on the validator.} The same feature service is known to be
incomplete on its \emph{line} layers -- its 400~kV layer holds only the western
leg of Moneypoint--Oldstreet, and no circuit at all for the legs east of
Oldstreet. \file{eirgrid.py} prints that warning. The station layer used here
shows no such gap, and the cross-check is station-to-station, but a validator
with a documented hole in it should be named as one.


Of the 134 cross-checked stations, one is more than 5~km out: TANDRAGEE, at
36.3~km. It is EirGrid's register that is wrong -- it puts the station at the
Republic's end of the Louth--Tandragee 275~kV interconnector, in Co.~Louth,
while the substation itself is in Co.~Armagh where OSM has it.

The 2033 cases place 78\% rather than 87\%. The difference is exactly what the
forecast statement is for: stations that do not exist yet are not in
OpenStreetMap.

% ===================================================================
\section{The North-West region}
\label{sec:northwest}

EirGrid and SONI run a Wind Dispatch Tool (WDT) in their control centres that
applies MW limits to individual wind and solar farms, or to groups of them
associated with a particular network constraint; the group definitions are
published \citep{eirgrid_wdt2024}. The region studied here -- Donegal, Sligo
and north Mayo -- is a heavily wind-dominated area behind a narrow export path,
and the brief for this work identified it as WDT constraint groups 1--3.
\file{northwest.py} extracts it and reconciles it against a hand-built 15-node
dataset.

\textbf{What the published groups actually are, and why it matters.} The
February 2024 overview lists the North West groups as: 1, Letterkenny A1 busbar
flows; 2, Letterkenny A2 busbar flows; 4, Carrick-on-Shannon to Arva; 5
(outage-driven), Sliabh Bawn to Lanesboro; 6 (outage-driven), Binbane to
Cathaleen's Fall; and 7, Corderry to Srananagh, retired after that circuit was
uprated. The 2020 version numbered them differently again -- group 1 was
Letterkenny to Clogher, group 2 Corderry to Srananagh
\citep{eirgrid_wdt2020,eirgrid_wdt2024}. Two consequences. The station set
studied here is clearly the right region: Letterkenny, Clogher, Binbane,
Cathaleen's Fall, Corderry and Srananagh are all in it, and all appear in
EirGrid's own group definitions. But the \emph{numbering} is not stable across
versions, so ``groups 1--3'' does not by itself pin down a fixed set of
stations, and nothing in this report should be read as reproducing a particular
numbered group.

I could not retrieve either overview document from this environment -- the
hosts are refused by its egress policy -- so the group names above come from a
search index of them rather than from the documents themselves. They should be
checked against the PDFs before being quoted onward.

\begin{figure}[H]\centering
\includegraphics[width=0.82\textwidth]{i_northwest.png}
\caption{The region as the kit ships it. Dot area is plant capacity, the ring
is demand, the square is Srananagh 220~kV -- the boundary and the reference.
1{,}329~MW of plant against 239~MW of demand: it exports, and the question is
whether the wires can carry it.}
\label{fig:northwest}
\end{figure}

\subsection{Two views, and the mapping between them}

The hand-built dataset folds generation into parent substations; TYTFS has
Golagh, Mulreavy, Cliff, Meentycat and Lenalea as their own 110~kV buses. Both
views are produced, and the mapping is explicit in one table in the code:

\begin{lstlisting}
STATIONS = (
    Station("Srananagh 220", (5042,), 220.0, note="the region's export path"),
    Station("Srananagh 110", (5041,), 110.0, parent="Srananagh 220"),
    Station("Cathaleen's Fall", (1701, 17010, 17061), 110.0),
    Station("Cliff", (1761,), 110.0, parent="Cathaleen's Fall"),
    ...
)
\end{lstlisting}

\where{northwest.py}{STATIONS, station\_table(), bus\_map(), extract()}

\subsection{The circuit count}

The brief said 16 circuits among the 15 stations. Reproduced independently:
\textbf{16 routes, carried by 19 circuits}. Both numbers are right and they
answer different questions -- a route is a pair of stations, a circuit is a
wire, and some routes are double-circuit. The count agrees.

\subsection{What folding to 15 nodes costs}

Same impedances, same injections, a different reference bus: a DC flow has to
put the same MW on every circuit as the full network does. That is a check with
a number attached.

\begin{table}[H]\centering\small
\begin{tabular}{lr}
\toprule
\textbf{what is folded} & \textbf{worst disagreement vs.\ the full network} \\
\midrule
nothing (24-station native view)   & \textbf{0.0099 MW} \\
all nine folds (15-station view)   & 21.59 MW \\
all but Srananagh 110              & \textbf{1.74 MW} \\
all but Srananagh 110, Lenalea, Cunghill & 0.15 MW \\
\bottomrule
\end{tabular}
\caption{Almost all of the error is one fold. Merging Srananagh's two busbars
short-circuits the 250~MVA transformer, whose 0.064~pu of reactance is a large
fraction of the region's total.}
\end{table}

\textbf{The recommendation is one node.} Split Srananagh back into 110 and 220
with the transformer between them -- 16 nodes instead of 15 -- and the flow
error falls from 21.6~MW to 1.7. That is a concrete, cheap change to a dataset
that is otherwise sound.

\where{northwest.py}{compare\_with\_full(), extract(), verify()}

\subsection{Three findings about the hand-built data}

\begin{enumerate}
\item \textbf{The hydro capacities are all identically 38.} Cathaleen's Fall,
      Cliff and Golagh have very different real capacities -- Cathaleen's Fall
      alone is two machines of 22.5 and 23~MW in the TYTFS file. 38 is the
      38~kV connection voltage entered into a capacity column: 38~kV is a real
      level in this system (278 buses carry it in WP2024), and Cathaleen's Fall
      has a bus named \code{HY\_CATH FALL} sitting at exactly 38~kV.
\item \textbf{The 24-node and 15-node blocks disagree by exactly Cunghill's
      35~MW.} One block includes it, the other does not.
\item \textbf{The supply set is a correct subset, not a wrong dataset.} Every
      plant in it is in TYTFS at the station it says. It is missing plant TYTFS
      has, which is a scope difference rather than an error.
\end{enumerate}

\subsection{Where the constraint actually is}

\code{python northwest.py windy --capacity-factor X} puts the region's
capacity to work at $X$ and re-solves \emph{the whole transmission network} --
not the extract, because an extract with five of six ties pinned sends every
extra megawatt out through the sixth and gives an answer that is both large and
meaningless.

\begin{table}[H]\centering\small
\begin{tabular}{lrlrr}
\toprule
\textbf{capacity factor} & \textbf{net export} & \textbf{worst boundary tie} &
\textbf{Srananagh} & \textbf{worst internal} \\
\midrule
10\%  & $-97$ MW (importing) & Letterkenny--Strabane 0.36$\times$ & $-85$ MW & 0.40$\times$ \\
20\%  & 18 MW                & Letterkenny--Strabane 0.71$\times$ & $-62$ MW & 0.48$\times$ \\
\rowcolor{surface}
\textbf{30\%} & \textbf{134 MW} & \textbf{Letterkenny--Strabane 1.05$\times$} & $-39$ MW & 0.56$\times$ \\
60\%  & 480 MW               & Letterkenny--Strabane 2.09$\times$ & $+31$ MW & 0.80$\times$ \\
100\% & 942 MW               & Letterkenny--Strabane 3.48$\times$ & $+123$ MW & \textbf{1.11$\times$} \\
\bottomrule
\end{tabular}
\caption{The binding constraint is the Letterkenny--Strabane 110~kV tie into
Northern Ireland -- \code{3581-89516-1}, from Letterkenny (area 1, IE) to the
Strabane phase-shifter bus (area 14, NI) -- and it binds at about 29\% of
connected capacity. Srananagh does not even reverse into export until about
45\%. \textbf{The rating is scenario-dependent}: \code{RATE1} on this circuit
is 93~MVA in WP2024, the case this sweep runs on, but 80 in SV2024, 105 in
SV2033 and 123 in WP2033. Quoting one number for it across scenarios, as the
first draft of this report did, is wrong.}
\label{tab:windy}
\end{table}

So Srananagh 220 is the right place to \emph{slack} the extraction and the
wrong place to look for the constraint. Remember this table; Section~\ref{sec:crossval}
comes back to it.

\where{northwest.py}{windy(), boundary\_flows(), region\_flows()}

\subsection{Both north-western cross-border ties have a phase shifter on them}
\label{sec:pst}

This came out of chasing a numerical discrepancy (Section~\ref{sec:kit}) and is
worth stating on its own. Reading the transformer records directly:

\begin{table}[H]\centering\small
\begin{tabular}{lllr}
\toprule
\textbf{from} & \textbf{to} & \textbf{rating} & \textbf{shift} \\
\midrule
89510 STRABANE (NI) & 89515 STRA\_PST & 125 MVA & $17^\circ$ \\
79010 ENNISKIL (NI) & 79015 ENNK\_PST & 125 MVA & $3^\circ$ \\
\bottomrule
\end{tabular}
\caption{The only two phase-shifting transformers in the transmission network,
in every one of the four cases. \code{PST} is the file's own naming.}
\end{table}

The circuit this report keeps returning to, \code{3581-89516-1}, runs from
Letterkenny straight to \code{89516 STRA\_PST} -- so the binding element sits
directly behind a phase shifter, which is the device a system operator installs
precisely to control flow on a tie. The other north-western cross-border tie,
Corraclassy to \code{79016 ENNK\_PST}, has one too.

\textbf{This model holds both taps fixed.} That is a sharper and more honest
statement of the limitation than the first draft's ``SONI would redispatch
around it'': there is a specific, identified control device on the constrained
path whose action is not modelled here.

\where{psse.py}{read\_raw() $\rightarrow$ case.transformer, field ANG1}

\subsection{The line lengths are real, and the first draft said otherwise}
\label{sec:lengths}

The first draft listed ``TYTFS line lengths are placeholders'' as a limitation.
Checking it rather than repeating it:

\begin{itemize}
\item \textbf{83\% of transmission branches carry a non-zero length} (561 of
      679 in WP2024), median 11.7~km, longest 208~km.
\item Of the 118 that do not, \textbf{111 are couplers and stubs} with
      $|X| \leq 0.01$, where a zero length is the correct value. Only seven are
      real circuits missing a length, and six of those seven are in Northern
      Ireland.
\end{itemize}

\begin{figure}[H]\centering
\includegraphics[width=\textwidth]{j_lengths.png}
\caption{Left: for the 420 circuits with a stated length and both ends
independently geocoded, the stated length against the great-circle distance
between the ends. $r = 0.983$, median ratio 1.17 -- a real route is a little
longer than a straight line, which is exactly what this shows. Two datasets
that share no source agreeing like this checks the geocoding as much as it
checks the lengths. Right: reactance per km of stated length at 110~kV.}
\label{fig:lengths}
\end{figure}

The right-hand panel is bimodal, and the split is physical rather than an
artefact. The 347 circuits above 0.25~$\Omega$/km have a median of
\textbf{0.394}~$\Omega$/km and a median shunt susceptance of
0.00035~pu/km; the 117 below have a median of 0.13~$\Omega$/km and
\textbf{0.011}~pu/km -- thirty times the charging. High capacitance and low
reactance per km is the signature of underground cable, and the names in that
cluster are Ringsend--Whitebank, Galway--Salthill, Marina--Trabeg: urban and
coastal, where cable is what gets laid.

Two things follow. The lengths are usable, with the seven gaps named. And
\code{gridkit.OHMS\_PER\_KM = 0.4}, which the first draft called a rule of
thumb, is within 2\% of the median of the overhead-line circuits in this very
dataset -- so a new line costed off it is being costed off the file's own
overhead-line population, which is worth saying instead.

% ===================================================================
\section{Time series: what is real and what is not}
\label{sec:profiles}

TYTFS gives four snapshots. A constraint study needs hours.

\subsection{The real pipeline, and why there is no real data}

\file{profiles.py} is a complete pipeline against Open-Meteo's Historical
Weather API \citep{openmeteo}, which serves ECMWF's ERA5 reanalysis
\citep{era5} on a 0.25\textdegree{} grid and the higher-resolution ERA5-Land
\citep{era5land} on a 0.1\textdegree{} grid, without a key and under CC BY 4.0.
It does batched requests,
disk caching by cell-year so nothing is ever fetched twice, a stated power
curve smoothed into a farm curve by Gaussian convolution after
\citet{staffell2016}, a hub-height correction from the two ERA5 wind levels, a
panel model using the Erbs diffuse-fraction correlation
\citep{erbs1982}, and a validation step against EirGrid's published all-island
wind series.

\textbf{It has never run, because both hosts are refused by this environment's
egress policy} -- 403 at CONNECT, before any request is made:

\begin{lstlisting}[language=shell]
archive-api.open-meteo.com:443   gateway answered 403 to CONNECT
www.smartgriddashboard.com:443   gateway answered 403 to CONNECT
cms.eirgrid.ie:443               gateway answered 403 to CONNECT
\end{lstlisting}

A 403 at CONNECT is an organisation policy denial. It was not retried and no
mirror was sought.

\textbf{So there are no ERA5 profiles in this repository and no validation
result.} What there is instead is the pipeline, tested, with the fetch as the
only untested step -- and every path that cannot get real data raises
\code{profiles.Unavailable} rather than falling back to something invented.
That is deliberate, and it is tested:

\begin{lstlisting}
test_loading_a_cell_that_was_never_fetched_raises
test_a_partial_year_is_not_cached_as_a_complete_one
test_a_response_missing_a_variable_is_refused
test_an_unreadable_eirgrid_file_raises_rather_than_guessing
\end{lstlisting}

The brief was explicit that the validation is the only thing standing between a
real profile and a plausible-looking fabrication. Shipping a synthetic year
under the ERA5 label, however carefully footnoted, would be that fabrication.

\textbf{Two ways to finish it.} Allowlist the two hosts and re-run -- total
traffic is about six HTTP requests per case-year. Or drop a hand-downloaded
dashboard CSV into \file{data/raw/eirgrid/} and the Open-Meteo JSON into
\file{data/raw/weather/}: everything after the fetch runs unchanged.

\where{profiles.py}{fetch(), load\_cell(), farm\_curve(), solar\_capacity\_factor(), validate\_wind(), Unavailable}

\subsection{The synthetic generator, and the one detail that matters}

\file{synthetic.py} generates profiles instead. It exists for two reasons: as
insurance against exactly the blocker above, and so participants can build
counterfactuals no historical year contains.

The single thing a hand-rolled generator gets wrong is spatial correlation.
Independent per-bus noise makes curtailment behaviour meaningless: 400 sites of
independent noise give a fleet output that barely leaves its mean, and a
network is never asked to carry anything. Real calm spells and real storm
fronts move across a country and hit sites together.

So the wind field is built as a Gaussian random field with an exponential
correlation decay in distance and an Ornstein--Uhlenbeck process in time, then
mapped through a Gaussian copula to a Weibull distribution:

\[
\rho(d) = \exp\!\left(-\frac{d}{L}\right), \qquad L = 400~\mathrm{km}
\]

\begin{figure}[H]\centering
\includegraphics[width=0.85\textwidth]{h_correlation.png}
\caption{The check that matters, and it is measured \emph{back off the output}
rather than assumed. Every pair of wind sites, correlation of hourly capacity
factor against distance. The black curve is the field it was built from; the
orange line is what came out. They agree. Had this come out flat and near zero,
every curtailment number in this report would be worthless.}
\label{fig:correlation}
\end{figure}

\where{synthetic.py}{correlation\_matrix(), gaussian\_field(), weibull\_from\_normal(), spatial\_check()}

\subsection{Anchored on the TYTFS states}

The generated year is not free-floating. WP2024 and WP2033 are winter-peak
conditions, SV2024 and SV2033 summer-valley, so the generator places each case
at an hour of the year that reproduces its state and tapers back to the
unconstrained field over 18 hours either side.

\begin{lstlisting}
ANCHOR_POINTS = {"WP": (1, 17, 18), "SV": (7, 9, 5)}   # month, day, hour
ANCHOR_WINDOW_H = 18
SEED = 42
\end{lstlisting}

Deterministic, seeded with 42, matching the convention in the hand-built
dataset's own README.

\where{synthetic.py}{anchors(), anchor\_index(), taper(), build()}

\subsection{Does WP2033 actually bind?}

The exercise turns on it. If the profiles do not produce hours where the
network cannot take what the wind offers, the central problem does not appear.

\begin{table}[H]\centering\small
\begin{tabular}{lrr}
\toprule
& \textbf{WP2024} & \textbf{WP2033} \\
\midrule
hours studied (windiest) & 60 & 60 \\
energy offered   & 1.4 GWh & 342.8 GWh \\
energy curtailed & \textbf{0.0 GWh} & \textbf{20.4 GWh} \\
curtailment      & 0.0\% & \textbf{5.94\%} \\
hours with curtailment & 0 of 60 & \textbf{60 of 60} \\
\bottomrule
\end{tabular}
\caption{It binds, in every hour studied. The 2024 contrast is not a fair fight
and should not be read as one -- that case runs almost no wind, so there are
only 1.4~GWh to curtail. What it shows is that the \emph{scenario}, not the
profile generator, creates the problem.}
\end{table}

\where{synthetic.py}{binding(), fleet\_variability()}

% ===================================================================
\section{The participant kit}
\label{sec:kit}

\file{participant-kit/} is a standalone directory. It contains no build
tooling and imports nothing from the rest of the repository.

\begin{table}[H]\centering\small
\begin{tabular}{ll}
\toprule
\file{networks/} & 4 scenarios $\times$ 2 scopes, as \code{.nc} and as CSV folders \\
\file{gridkit.py} & load, edit, reset, save; and read results back \\
\file{flowmath.py} & PTDF, susceptibility, shift factors \\
\file{plotstyle.py} & one validated palette for every chart \\
\file{examples/} & six scripts, a to f \\
\file{test\_kit.py} & 20 checks, runnable without pytest \\
\bottomrule
\end{tabular}
\end{table}

Every network carries 168 hourly snapshots -- one week, centred on the hour
that matches its own TYTFS state. Both formats ship because many participants
would rather edit a spreadsheet than learn an API.

\subsection{The loader}

\begin{lstlisting}
import gridkit

gridkit.quiet()                                  # PyPSA and HiGHS are chatty
n = gridkit.load("WP2033", "all-island")         # or "north-west"

gridkit.add_line(n, "Meentycat", "Srananagh 220", s_nom=200)
gridkit.remove_line(n, "3581-89516-1")
gridkit.set_rating(n, "2781-4951-1", 200.0)
gridkit.add_battery(n, "Letterkenny", p_nom=50, hours=4)

gridkit.solve(n)                                 # n.optimize, log off
print(gridkit.curtailment(n).head())
print(gridkit.binding(n).head())

n = gridkit.reset(n)                             # back to the baseline
\end{lstlisting}

\where{participant-kit/gridkit.py}{load(), add\_line(), remove\_line(), set\_rating(), add\_battery(), reset(), curtailment(), binding(), unserved()}

\subsection{Two traps, both already sprung}

These are in the kit's README because both cost time to find and neither
announces itself.

\textbf{\code{freeze\_dispatch}.} \code{n.optimize()} writes its answer to
\code{generators\_t.p}. \code{n.lpf()} reads \code{generators\_t.p\_set}, which
is a different thing and which the kit ships empty. Run \code{lpf} straight
after \code{optimize} and you get a power flow with no generation in it: the
reference bus silently supplies the whole island, and one circuit shows at
1600\% of its rating.

\textbf{Buses at 0$^\circ$N 0$^\circ$E.} PyPSA's netCDF writer turns a missing
coordinate into \code{0.0} rather than \code{NaN}, and 0$^\circ$N 0$^\circ$E is
in the Gulf of Guinea. Plot the raw \code{n.buses} and Ireland is squashed into
the corner of an Atlantic-to-equator map. \code{gridkit.placed\_buses(n)}
drops those and respects the \code{coordinate\_source} column.

\clearpage
\subsection{(a) DC power flow}

\begin{figure}[H]\centering
\includegraphics[width=0.72\textwidth]{a_flow_map_WP2033_all-island.png}
\caption{WP2033 at peak demand. Each circuit is coloured by how full it is; the
one at or over its rating is drawn in red and ringed, because several circuits
join buses a few hundred metres apart and a stroke alone would be invisible.
Nothing here is optimised -- a DC power flow takes the dispatch as given and
tells you where the flow ends up.}
\end{figure}

Six circuits sit at their rating for at least one hour of the week. The worst,
\code{3581-89516-1}, is at 100\% for \textbf{118 of 168 hours}.

\where{participant-kit/examples/a\_dc\_power\_flow.py}{main(), \_draw()}

\clearpage
\subsection{(b) Dispatch and curtailment -- and how much of it is the network}

\begin{figure}[H]\centering
\includegraphics[width=\textwidth]{b_dispatch_WP2033_all-island.png}
\caption{Least-cost dispatch over the week, and the reason for the
curtailment.}
\label{fig:dispatch}
\end{figure}

Wind supplies 68\% of the week's energy. 291~GWh of renewable output is
curtailed -- which sounds like a damning verdict on the transmission system
until you separate it.

\textbf{The separation is one extra solve.} Lift every branch rating out of the
way and re-optimise. Whatever is \emph{still} curtailed is surplus that no
network could have taken; the difference is what the transmission cost.

\begin{table}[H]\centering\small
\begin{tabular}{lrl}
\toprule
265.2 GWh & 91\% & surplus -- 42.6~GW of plant against an 8.8~GW peak \\
\rowcolor{surface}
\textbf{26.3 GWh} & \textbf{9\%} & \textbf{congestion -- what the network refused} \\
\bottomrule
\end{tabular}
\caption{Curtailment on its own does not mean the network is the problem. In
WP2033 most of it is not, and building a line recovers none of the 265~GWh.
This distinction is the difference between a headline and a finding.}
\end{table}

\where{participant-kit/examples/b\_lopf\_dispatch.py}{\_congestion\_share(), by\_carrier()}

\clearpage
\subsection{(c) Capacity expansion}

\begin{figure}[H]\centering
\includegraphics[width=0.92\textwidth]{c_expansion_WP2033_north-west.png}
\caption{The binding circuits are marked extendable, a battery is offered at
every renewable bus, both get an annualised cost, and the optimiser trades
building against curtailing.}
\end{figure}

On the North-West region it reinforces one circuit by 83~MVA and builds three
batteries, cutting the week's curtailment from 3.7~GWh to 0.6. The costs are
round numbers rather than a price list -- the point is the pattern: mark a
component extendable, give it a capital cost, solve, read \code{s\_nom\_opt}
back.

\where{participant-kit/examples/c\_capacity\_expansion.py}{\_make\_extendable(), \_offer\_batteries()}

\clearpage
\subsection{(d) PTDF from the pseudoinverse}

Everything in \file{flowmath.py} comes out of one object, the weighted graph
Laplacian $L = K B K^{\!\top}$, where $K$ is the bus--branch incidence matrix
and $B$ the diagonal of branch susceptances.

$L$ is singular -- adding a constant to every angle changes no flow -- so it is
inverted with the Moore--Penrose pseudoinverse:
\[
\theta = L^{+}(p + KB\varphi), \qquad
F = B\,(K^{\!\top}\theta - \varphi), \qquad
\mathrm{PTDF} = B K^{\!\top} L^{+}
\]

\textbf{Why the pseudoinverse rather than deleting a slack row.} The textbook
fix is to pick a slack, delete its row and column, and invert what is left.
That works, and it writes the arbitrary choice of slack into every number that
comes out. $L^{+}$ instead returns the solution orthogonal to the nullspace --
the one where the balancing megawatt is spread evenly. The choice has not
disappeared; it has become explicit, and \code{shift\_factors()} lets you
change it without rebuilding anything.

\begin{figure}[H]\centering
\includegraphics[width=\textwidth]{d_ptdf_WP2033_all-island.png}
\caption{Left: one PTDF column on the map -- how much of an injection at each
bus the monitored circuit carries. The pale middle of the island is the part
the circuit cannot feel at all. Right: the check. Branch flows rebuilt from
$\mathrm{PTDF}\times p$ against what \code{n.lpf()} produced, over 980
branches. Maximum disagreement $5.5\times10^{-10}$~MW.}
\label{fig:ptdf}
\end{figure}

\textbf{That check earned its keep.} It first came out at \textbf{79~MW}, and
the reason was the network's two phase-shifting transformers -- one of them at
17$^\circ$. A phase shifter imposes an angle of its own, so the flow is
$b(\theta_i - \theta_j - \varphi)$ and the bus equation gains a $KB\varphi$
term; it behaves exactly like a pair of equal and opposite injections. The
PTDF is unaffected, the flows are not. A PTDF that has not been checked against
a power flow is a matrix of plausible numbers.

\where{participant-kit/flowmath.py}{branches(), laplacian(), pseudoinverse(), ptdf(), angles(), flows()}

\clearpage
\subsection{(e) Edge-to-edge susceptibility, and Braess's paradox}
\label{sec:braess}

Differentiating $F = BK^{\!\top}L^{+}p$ with respect to one branch's
susceptance gives the full edge-to-edge response matrix:
\[
\frac{\partial F_e}{\partial B_{e'}}
= \delta_{ee'}\,a_e \;-\; B_e\left(k_e^{\!\top} L^{+} k_{e'}\right) a_{e'}
\]
where $a_e = F_e/B_e$ is the effective angle difference across branch $e$
(the plain angle difference, less any phase shift).

The diagonal is the obvious part. The off-diagonal is where \textbf{Braess's
paradox} lives \citep{braess1968}: a positive entry means that strengthening
$e'$ puts \emph{more} power on $e$. The effect is established in power networks
specifically -- theoretically by \citet{witthaut2012}, and demonstrated on a
real AC laboratory grid with a topological predictor by \citet{schafer2022}.
Power in a network is not routed, it is imposed --
it divides between parallel paths in inverse proportion to reactance, and
nothing consults the ratings. Build a second circuit anywhere in a mesh and
every impedance ratio in it changes.

\begin{figure}[H]\centering
\includegraphics[width=\textwidth]{e_braess_WP2033_all-island.png}
\caption{Left: the branches that move the island's most-loaded circuit, signed.
Of the 979 other branches, strengthening \textbf{516} of them makes it worse.
Right: the derivative is a claim, so the script tests it. It doubles the
susceptance of the strongest candidate -- what a second circuit on the same
route does -- holds the dispatch fixed, and re-solves the flow.}
\label{fig:braess}
\end{figure}

\textbf{The result.} Doubling \code{3581-28019-1} makes \code{3581-89516-1}
carry more power in \textbf{168 of 168 hours}. Flow at the studied hour goes
from 123.0 to 127.3~MW, against a 123~MVA rating. A reinforcement that a
reasonable planner might propose makes the binding constraint worse.

\where{participant-kit/flowmath.py}{susceptibility(), braess\_candidates()}

\clearpage
\subsection{(f) Shift factors: the Wind Dispatch Tool calculation}

For one monitored circuit, how much of each generator's output that circuit
carries. Rank the generators by it and you have drawn something with the shape
of a constraint group: the machines a dispatch instruction would have to act
on.

\textbf{``Reproducing EirGrid's WDT methodology'' overstates this, and the
heading is the brief's wording rather than a claim.} EirGrid's published groups
are defined by a monitored flow and are the product of contingency studies, an
AC network, and the Northern Ireland system \citep{eirgrid_wdt2024}; what
follows is a DC shift-factor ranking on one intact circuit. It is the right
arithmetic and the wrong amount of it.

\begin{figure}[H]\centering
\includegraphics[width=\textwidth]{f_shift_factors_WP2033_all-island.png}
\caption{Left: the shift factor at every bus, dot area being installed
capacity. Right: the constraint group, ranked by the MW of relief each member
could actually deliver at the studied hour.}
\end{figure}

The one thing to be careful about is the \emph{reference}. A shift factor is
``MW on the circuit per MW at the generator'', but a megawatt cannot appear on
its own, so the definition is incomplete until you say where the balancing
megawatt goes. The kit offers three conventions:

\begin{table}[H]\centering\small
\begin{tabular}{ll}
\toprule
\code{reference="load"} & spread over every load in proportion to its size --
what a system operator means \\
\code{reference="uniform"} & spread evenly over all buses -- the
pseudoinverse's own reference \\
a bus name & all of it at that bus -- the textbook single-slack convention \\
\bottomrule
\end{tabular}
\end{table}

\textbf{The values differ; the ranking does not.} Spearman rank correlation
between the three references is \textbf{1.0000}. That invariance is why the
method is robust enough to run a real dispatch tool on -- and it is worth
knowing that it is an invariance rather than a coincidence.

\where{participant-kit/flowmath.py}{shift\_factors()}

% ===================================================================
\clearpage
\section{Findings}

\subsection{Two independent methods agree on the binding circuit}
\label{sec:crossval}

This is the result worth the most, because nothing was tuned to produce it.

\begin{table}[H]\centering\small
\begin{tabularx}{\textwidth}{lXX}
\toprule
& \textbf{Section~\ref{sec:northwest}} & \textbf{Section~\ref{sec:kit}} \\
\midrule
network   & 2024 & 2033 \\
method    & DC power flow, other ties pinned at the case's own values &
            full LOPF over a week \\
profiles  & none -- a capacity-factor sweep & synthetic correlated weather \\
scope     & the North-West region & the whole island \\
\midrule
rating of that circuit & 93 MVA (\code{RATE1} in WP2024) & 123 MVA
            (\code{RATE1} in WP2033) \\
\midrule
\textbf{answer} & \multicolumn{2}{c}{\cellcolor{surface}\textbf{\code{3581-89516-1},
Letterkenny to the Strabane phase-shifter bus}} \\
\bottomrule
\end{tabularx}
\caption{Different network vintage, different method, different data, no shared
code beyond the conversion itself. Same circuit. A single method that finds a
constraint has found a property of its own assumptions until something else
agrees with it.}
\end{table}

\textbf{How much this is worth, stated carefully.} The two methods share the
conversion, the geocoding and the topology; they are independent in vintage,
dispatch, profiles and scope, not in everything. And the agreement is partly
structural: the region exports, and this is its narrowest export path, so a
thermal-limit method was always likely to land here. What the agreement rules
out is that the finding is an artefact of the synthetic weather or of the
2033 build-out -- not that it is the constraint a control room would act on.

Independent support from outside this work: EirGrid's two live North West
constraint groups in February 2024 both monitor \emph{Letterkenny busbar flows}
\citep{eirgrid_wdt2024}, and the circuit found here is one of the circuits on
that busbar. That is corroboration of the area, not of the specific circuit,
and Section~\ref{sec:northwest} explains why it cannot be more than that.

\subsection{Most curtailment is not the network's fault}

Section~\ref{sec:kit} again: of 291~GWh curtailed in WP2033, 265 is surplus and
26 is congestion. The instinct on seeing a large curtailment figure is to
propose reinforcement. In this scenario, 91\% of it would be money spent on
energy that has nowhere to go regardless.

\subsection{Reinforcement can make it worse, and here is a case}

Section~\ref{sec:braess}. Of the 979 branches that could be strengthened, 516
would increase the flow on the island's most-loaded circuit. The strongest
candidate, verified on a re-solved power flow, makes it worse in every hour of
the week. This is not a curiosity -- it is a reason to compute the
susceptibility matrix before proposing a circuit.

\subsection{Three findings about the hand-built dataset}

Restated from Section~\ref{sec:northwest} because they are actionable: the
hydro capacities are the 38~kV connection voltage rather than capacities; the
24-node and 15-node blocks differ by exactly Cunghill's 35~MW; and splitting
Srananagh back into two nodes takes the model's flow error from 21.6~MW to 1.7.

% ===================================================================
\section{Bugs found, and what each cost}

Each of these produced a network that ran. That is the reason to write them
down.

\begin{enumerate}[leftmargin=1.5em]

\item \textbf{Every LOPF infeasible.} PyPSA 1.x turns a non-null \code{p\_set}
into a hard equality constraint pinning the dispatch. The networks carried the
TYTFS snapshot's own dispatch, which is the case's answer asserted -- and
infeasible against a lossless model. Fixed by
\code{pypsa\_net.for\_optimisation()}, which clears it and keeps the case's
values in a column of their own.

\item \textbf{The full network still infeasible.} The coupler-bounding rule,
applied to sub-110~kV distribution transformers, produced limits that bind on
the case's own flows -- 13 transformers, worst at 156\%. Fixed by restricting
the rule to zero-impedance couplers.

\item \textbf{WP2033's angles out by 100$^\circ$.} PyPSA's own slack choice
survived a per-label assignment; one 21~MW hydro machine absorbed 856~MW.
Section~\ref{sec:slack}. Angle correlation 0.90 $\rightarrow$ 0.99.

\item \textbf{The LOPF exported 1{,}223~MW for profit.} A single bidirectional
generator with a positive marginal cost and \code{p\_min\_pu = -1} gets
\emph{paid} to run backwards. Split into an import generator and a
\code{sign=-1} export sink.

\item \textbf{7.8~GW of ``wind'' running flat out.} 82 wind and solar
generators in WP2033 had no \code{p\_max\_pu} series, because the synthetic
weather is generated at geocoded sites and their buses never matched
OpenStreetMap. PyPSA defaults a missing series to 1.0, so at a zero bid they
ran at full output for all 168 hours -- more than the island's peak demand, a
fifth of the fleet, must-run baseload labelled ``wind''. The curtailment figure
read 64\% instead of 26\%. Each now borrows the nearest profile of its own
carrier. \textbf{The dispatch chart looked entirely reasonable throughout.}

\item \textbf{The North-West had no wind at all.} \code{northwest.extract()}
folds every machine at a station into one generator with carrier \code{AC},
which was right for the reconciliation it was written for. For the kit it meant
a region that is 90\% wind could not curtail. Now split by carrier from the
case's own machine records.

\item \textbf{Flows out by 79~MW.} The two phase-shifting transformers,
Section~\ref{sec:kit}. Caught only because the PTDF was checked against a power
flow.

\item \textbf{A circuit at 1600\% of its rating.} \code{lpf} after
\code{optimize} without freezing the dispatch. Now
\code{gridkit.freeze\_dispatch()}, with the symptom described in its docstring
so the next person recognises it.

\item \textbf{Ireland in the corner of a map of the Atlantic.} Three buses at
0$^\circ$N 0$^\circ$E. Now \code{gridkit.placed\_buses()}.

\item \textbf{Guards that never fired.} Several \code{if not len(frame\_t.p)}
checks were testing the row count of an empty-column DataFrame -- which is the
snapshot count, and always truthy -- so \code{freeze\_dispatch} would silently
freeze nothing.

\end{enumerate}

% ===================================================================
\section{Tests}

\begin{lstlisting}[language=shell]
python -m pytest            # 215 passed
cd participant-kit && python test_kit.py    # 20/20, no pytest needed
\end{lstlisting}

\begin{table}[H]\centering\small
\begin{tabular}{lrl}
\toprule
\textbf{file} & \textbf{tests} & \textbf{what it guards} \\
\midrule
\file{test\_network.py}   & 37 & the earlier graph work \\
\file{test\_profiles.py}  & 32 & the ERA5 pipeline, and that it refuses to invent \\
\file{test\_synthetic.py} & 31 & correlation, anchoring, determinism \\
\file{test\_psse.py}      & 30 & the raw-file parser, against the files themselves \\
\file{test\_pypsa\_net.py} & 26 & every conversion decision in Section~\ref{sec:conversion} \\
\file{test\_northwest.py} & 21 & the station map, the circuit count, both views \\
\file{participant-kit/test\_kit.py} & 20 & the kit: networks, loader, linear algebra \\
\file{test\_geocode.py}   & 18 & normalisation, NI codes, the accept/reject rules \\
\midrule
& \textbf{215} & \\
\bottomrule
\end{tabular}
\end{table}

\subsection{The two that matter most}

Most tests check that code does what it says. These two check it against
something that does not share code with it.

\textbf{\code{test\_ptdf\_matches\_a\_power\_flow}} rebuilds every branch flow
from $\mathrm{PTDF}\times p$ and compares against PyPSA's own \code{lpf}. Two
routes to the same number, agreeing to $10^{-6}$~MW. This is the test that
found the phase shifters.

\textbf{\code{test\_susceptibility\_matches\_finite\_differences}} perturbs one
branch's susceptance by a ten-thousandth and checks the analytic derivative
against the difference it actually makes:

\begin{lstlisting}
numeric = (flows_with(bumped) - reference) / step
predicted = analytic.iloc[:, column].to_numpy()
# The natural scale is MW per MW/radian.  Some columns are exactly zero -
# a radial stub's susceptance changes no flow anywhere - so the column's
# own magnitude is not a usable denominator on its own.
scale = max(np.abs(predicted).max(), np.abs(reference).max() / base[column])
assert np.abs(numeric - predicted).max() / scale < 1e-4
\end{lstlisting}

If either drifts, nothing built on \file{flowmath.py} means anything.

\where{participant-kit/test\_kit.py}{test\_ptdf\_matches\_a\_power\_flow()}
\where{participant-kit/test\_kit.py}{test\_susceptibility\_matches\_finite\_differences()}

% ===================================================================
\section{Limitations}

Read this before presenting anything from it.

\textbf{This is a DC approximation of an AC model.} The TYTFS files are full AC
load-flow cases with reactive power, voltage magnitudes, tap changers and shunt
compensation. All of that is thrown away: flows are linear in angle
differences, voltages are 1.0~pu everywhere, losses are zero. That is the
standard approximation for transmission constraint work and is good to a few
percent on real flows, but it says nothing about voltage, reactive support or
stability -- and a conclusion that depends on those is not supported by
anything here.

\textbf{Line lengths are usable, with seven named gaps.} This corrects the
first draft, which called them placeholders. 83\% of transmission branches
carry a length; it correlates at $r = 0.983$ with the great-circle distance
between independently geocoded endpoints; and reactance per km separates
cleanly into overhead line and cable. Seven real circuits in WP2024 carry no
length, six of them in Northern Ireland, and a per-km calculation should skip
those rather than treat a zero as a distance. See Section~\ref{sec:lengths}.

\textbf{Ratings are TYTFS \code{RATE1}, not operational limits, and the reading
of \code{RATE1} is an inference.} See Section~\ref{sec:rate}: what the twelve
PSS/E rating slots mean is not publicly documented, and the normal/LTE/STE
reading here comes from the data's own structure. Beyond that, a control room
works to seasonal, weather-dependent and often dynamic limits that are not in
these files, and runs circuits above the continuous rating for defined periods.
``100\% of rating'' here means 100\% of a planning number, under an inferred
interpretation of which column that is. Ratings also change between scenarios:
the circuit this report follows is rated 80, 93, 105 and 123 MVA in the four
cases.

\textbf{Every time series is synthetic.} No hour in them ever happened. They
are anchored so each network's own scenario hour reproduces the TYTFS state,
and the spatial correlation is measured back off the output rather than
assumed -- but they are not a historical year and not a forecast. Generators
without a geocoded site borrow a neighbour's profile, which makes those pairs
perfectly correlated when they should not be.

\textbf{Costs are placeholders, and renewables bid negative.} Gas 90, biomass
40, hydro 1, imports 150 EUR/MWh, wind and solar at $-1$. The load-shedding
generators bid 10{,}000 EUR/MWh, which is a round number chosen to sit far above
everything else, \emph{not} the Single Electricity Market's Value of Lost Load
-- that was EUR 12{,}533/MWh for 2022 \citep{semc_voll2023}. The negative bid is
deliberate -- it is how a support scheme makes a wind farm willing to pay to
stay on, and it is what makes curtailment a last resort rather than a free
choice. It also makes the objective negative, so only \emph{differences}
between two objectives mean anything.

\textbf{About a fifth of generation carries the carrier \code{unknown}.}
Carriers are inferred from bus-name prefixes and a keyword list; the files do
not label plant by technology. Any result turning on the split between
\code{gas} and \code{unknown} is unsupported.

\textbf{The DC links are simplified.} Fixed efficiency, no ramp limits, no
minimum stable export, no market coupling. Interconnector flow here is what the
cost assumptions make of it, not what the day-ahead market would do.

\textbf{Only one circuit binds, the far side is an infinite sink, and there is
an unmodelled control device on the constrained path.} The model prices export
to Northern Ireland at nothing beyond the tie's thermal limit. In reality SONI,
the Northern Ireland transmission system operator \citep{soni}, would
redispatch around it -- and, more specifically, the tie terminates at a
$17^\circ$ phase-shifting transformer whose tap this model holds fixed
(Section~\ref{sec:pst}). What the result establishes is that the profiles
produce binding constraints and where, not that any particular curtailment
percentage is right.

\textbf{The TYTFS data is EirGrid's, provided for reference purposes only.} It
describes a planning forecast, not the operational system, and is published
without warranty as to accuracy or fitness for any particular purpose. Nothing
here is an EirGrid product or an EirGrid position.

% ===================================================================
\section{What this revision corrected, and what it could not check}
\label{sec:corrections}

The first version of this document was written from the working notes without
re-deriving the numbers or sourcing the background. Four things were wrong.

\begin{enumerate}[leftmargin=1.5em]

\item \textbf{``Line lengths are placeholders.'' False, and usefully so.} 83\%
of transmission branches carry a length and it correlates at $r = 0.983$ with
independently geocoded endpoint distances (Section~\ref{sec:lengths}). This
went from a limitation to a validation result -- and it incidentally validates
the geocoding a second time.

\item \textbf{``\code{RATE1 == RATE2} in all 1{,}129 branch records, in every
case.''} 1{,}129 is the WP2024 count alone; the four cases have 1{,}129,
1{,}120, 1{,}296 and 1{,}296, and the equality holds in all 4{,}841. The
related claim that \code{RATE3/RATE1} ``takes exactly two values'' also missed
97 summer-case records at 1.10021.

\item \textbf{``RATE1, RATE2 and RATE3 are normal / LTE / STE'' asserted as a
property of the PSS/E format.} It is an inference from this data. The format
carries twelve rating slots and Siemens' manual is not public
(Section~\ref{sec:rate}).

\item \textbf{One rating quoted for a circuit whose rating is
scenario-dependent.} \code{3581-89516-1} is 93 MVA in WP2024 and 123 in
WP2033, and the cross-validation table compared a WP2024 result with a WP2033
result under the WP2024 number.

\end{enumerate}

A fifth item is a reframing rather than an error: the report described the
study region as ``WDT constraint groups 1--3'' because the brief did. EirGrid's
published group numbering is not stable across versions, so that phrase does
not identify a fixed set of stations (Section~\ref{sec:northwest}).

\subsection{What could not be verified from here}

Stated plainly, because a citation the author could not open is weaker than one
they could.

\begin{itemize}
\item \textbf{Every EirGrid, SONI and SEM-O host is refused by this
      environment's egress policy}, as are \texttt{arxiv.org} and
      \texttt{nature.com}. So the TYTFS 2024 landing page, both Wind Dispatch
      Tool overviews and the SEM VoLL paper are cited as located publications;
      their contents reached this document through a search index, not through
      the documents. The WDT group names in Section~\ref{sec:northwest} are the
      claim most exposed to that, and they should be checked against the PDF.
\item \textbf{The PSS/E format itself.} Siemens' manual is not public. The
      bus \code{IDE} codes are corroborated by an independent open-source
      reader \citep{powsybl_psse}; the rating semantics are not corroborated by
      anything, which is why they are presented as an inference.
\item \textbf{Nothing in the TYTFS files was checked against the physical
      system.} Where this report says a circuit is rated at 93 MVA, it means
      the file says so.
\end{itemize}

Everything not in this list -- every table of measurements, every figure, every
count -- is produced by code in the repository and was re-derived for this
revision.

\section{Where everything is}

\begin{table}[H]\centering\small
\begin{tabular}{lll}
\toprule
\textbf{file} & \textbf{lines} & \textbf{what it does} \\
\midrule
\multicolumn{3}{l}{\color{inksoft}\textit{build half}} \\
\file{psse.py}        &  675 & reads PSS/E v35 raw files \\
\file{pypsa\_net.py}  & 1352 & the conversion, Section~\ref{sec:conversion} \\
\file{geocode.py}     & 1085 & OpenStreetMap matching, Section~\ref{sec:geocode} \\
\file{northwest.py}   &  958 & the region, Section~\ref{sec:northwest} \\
\file{profiles.py}    &  948 & the ERA5 pipeline (never run -- blocked host) \\
\file{synthetic.py}   &  964 & the synthetic generator, Section~\ref{sec:profiles} \\
\file{build\_kit.py}  &  459 & assembles the participant kit \\
\midrule
\multicolumn{3}{l}{\color{inksoft}\textit{participant half -- imports none of the above}} \\
\file{participant-kit/gridkit.py}   & 345 & the loader \\
\file{participant-kit/flowmath.py}  & 340 & PTDF, susceptibility, shift factors \\
\file{participant-kit/plotstyle.py} &  99 & the palette \\
\file{participant-kit/test\_kit.py} & 300 & the checks \\
\file{participant-kit/examples/}    & --- & six scripts, a to f \\
\midrule
\multicolumn{3}{l}{\color{inksoft}\textit{this document}} \\
\file{docs/report\_figures.py}  & 361 & the figures the kit does not draw \\
\file{docs/report/report.tex}   & --- & the source of this report \\
\file{docs/report/references.bib} & --- & 19 sources, each checked \\
\midrule
\multicolumn{3}{l}{\color{inksoft}\textit{the long-form working notes}} \\
\file{docs/PHASE2\_PYPSA.md}    & 605 & every conversion decision, with its evidence \\
\file{docs/PHASE3\_NORTHWEST.md} & 474 & the reconciliation, table by table \\
\file{docs/PHASE4\_PROFILES.md} & 338 & the ERA5 pipeline and the blocker \\
\file{docs/SYNTHETIC\_DATA.md}  & 333 & what is real and what is fabricated \\
\bottomrule
\end{tabular}
\caption{This report is the summary. The phase documents are the working notes,
and where a number here has a longer story, that is where it is.}
\end{table}

\vspace{5mm}
\noindent\textbf{Reproducing every figure in this document.}

\begin{lstlisting}[language=shell]
python docs/report_figures.py                 # the four repo figures
cd participant-kit
for e in examples/*.py; do python "$e"; done   # the six kit figures
cd ../docs/report && ./build.sh                # pdflatex, bibtex, pdflatex x2
\end{lstlisting}

\clearpage
\bibliography{references}

\end{document}
